\section*{Розділ II. Структура та склад СР СМ}
\addcontentsline{toc}{section}{Розділ II. Структура та склад СР СМ}

\subsection*{2.1. Структура СР СМ}
\addcontentsline{toc}{subsection}{2.1. Структура СР СМ}
    2.1.1. СР СМ є єдиним колегіальним органом, що здійснює свою діяльність через засідання її членів (далі – засідання СР СМ).

    2.1.2. Для забезпечення ефективного виконання своїх завдань та повноважень за окремими напрямами діяльності СР СМ формує зі свого складу або із залученням інших студентів постійні або тимчасові робочі органи (комітети, комісії, робочі групи тощо). Порядок їх формування та діяльності визначається цим Положенням та/або внутрішніми документами СР СМ.

\subsection*{2.2. Склад СР СМ}
\addcontentsline{toc}{subsection}{2.2. Склад СР СМ}
    2.2.1. Членами СР СМ з правом голосу є:

        \begin{enumerate}[label=\alph*)]
            \item Голова СР СМ, обраний на прямих виборах серед мешканців гуртожитків студмістечка;
            \item Голови Студентських рад гуртожитків (СРГ) за посадою (ex officio) на час виконання ними повноважень голів СРГ.
        \end{enumerate}

    2.2.2. СР СМ має право залучати до своєї роботи інших студентів Університету денної форми навчання для виконання певних функцій (наприклад, як керівників або членів комітетів, секретарів тощо). Такі залучені студенти беруть участь у роботі СР СМ з правом дорадчого голосу та не вважаються членами СР СМ у контексті прийняття рішень, що потребують голосування.

    2.2.3. Кількісний склад СР СМ з правом голосу дорівнює кількості гуртожитків у студмістечку Університету плюс одна особа (Голова СР СМ).

\subsection*{2.3. Вимоги до членів та кандидатів}
\addcontentsline{toc}{subsection}{2.3. Вимоги до членів та кандидатів}
    2.3.1. Кандидатом на посаду Голови СР СМ може бути будь-який студент Університету денної форми навчання, який відповідає вимогам, встановленим Положенням про Центральну виборчу комісію студентів (ЦВКс).

    2.3.2. Голова СР СМ не може одночасно обіймати посаду Голови Студентської ради гуртожитку.

    2.3.3. Членом СР СМ за посадою може бути виключно Голова СРГ, обраний відповідно до встановленого порядку.

    2.3.4. Студенти, які залучаються до роботи СР СМ з правом дорадчого голосу, повинні бути студентами Університету денної форми навчання.

\subsection*{2.4. Строк та порядок набуття і припинення повноважень}
\addcontentsline{toc}{subsection}{2.4. Строк та порядок набуття і припинення повноважень}
    2.4.1. Голова СР СМ обирається мешканцями гуртожитків студмістечка шляхом прямого таємного голосування строком на 1 (один) рік відповідно до процедури, визначеної Положенням про ЦВКс. Одна і та сама особа не може обіймати посаду Голови СР СМ більше двох строків поспіль.

    2.4.2. Повноваження Голови СР СМ можуть бути достроково припинені за рішенням Конференції студентів Університету (КСУ) у порядку, визначеному Положенням про ОСС та/або Регламентом КСУ, або за рішенням Студентської уповноваженої делегації (СУД) у випадках та порядку, передбачених Положенням про СУД.

    2.4.3. Строк повноважень членів СР СМ – Голів СРГ – становить 1 (один) рік та збігається зі строком їхніх повноважень на посаді Голови СРГ.

    2.4.4. Член СР СМ – Голова СРГ – набуває повноважень з моменту його обрання Головою СРГ та припиняє їх у разі припинення повноважень Голови СРГ.

    2.4.5. Повноваження студентів, залучених до роботи з правом дорадчого голосу, визначаються рішенням СР СМ та можуть бути припинені достроково за рішенням СР СМ, або у зв'язку із завершенням роботи відповідного робочого органу чи закінченням строку повноважень складу СР СМ, що їх залучив.
